\documentclass{article}
\usepackage{amsmath,amssymb,amsthm}
\usepackage{algorithm,algorithmic}
\usepackage{enumitem}
\usepackage{hyperref}
\usepackage{geometry}
\geometry{margin=1in}
\newtheorem{theorem}{Theorem}
\newtheorem{lemma}{Lemma}
\begin{document}

\section*{Algorithm Name}
\textbf{Randomized Block Coordinate Descent (RBCD) for Non‑Convex Smooth Functions}

\section*{Mathematical Formulation}
Let $w\in\mathbb{R}^{d}$ be the optimisation variable.  
Partition the coordinate set $\{1,\dots ,d\}$ into $B$ (disjoint) blocks  
\[
\mathcal{S}_{1},\dots ,\mathcal{S}_{B},\qquad 
\bigcup_{b=1}^{B}\mathcal{S}_{b}=\{1,\dots ,d\}.
\]
For a vector $w$, denote by $w^{(b)}\in\mathbb{R}^{|\mathcal{S}_{b}|}$ the sub‑vector restricted to block $b$, and by $\nabla_{b}f(w)$ the gradient of $f$ w.r.t.\ this block.

At iteration $t\in\{0,1,\dots ,T-1\}$:

\begin{enumerate}[label=\arabic*)]
\item Sample a block index $i_{t}\in\{1,\dots ,B\}$ independently and uniformly:
      \[
      \Pr(i_{t}=b)=\frac{1}{B},\qquad b=1,\dots ,B .
      \]
\item Compute the (exact) block gradient $\nabla_{i_{t}}f(w_{t})$.
\item Perform a gradient step on the selected block only:
\[
w_{t+1}^{(i_{t})}=w_{t}^{(i_{t})}-\eta_{t}\,\nabla_{i_{t}}f(w_{t}),\qquad 
w_{t+1}^{(b)}=w_{t}^{(b)}\;\;\forall b\neq i_{t}.
\]
\end{enumerate}
The step‑size sequence $\{\eta_{t}\}_{t\ge 0}$ is a scalar hyperparameter (possibly diminishing).

\section*{Pseudocode}
\begin{algorithm}[H]
\caption{Randomized Block Coordinate Descent (RBCD)}
\begin{algorithmic}[1]
\STATE \textbf{Input:} initial point $w_{0}\in\mathbb{R}^{d}$, step sizes $\{\eta_{t}\}$, number of iterations $T$.
\FOR{$t=0$ \TO $T-1$}
    \STATE Sample $i_{t}\sim \text{Uniform}\{1,\dots ,B\}$.
    \STATE Compute block gradient $g_{t}\gets \nabla_{i_{t}}f(w_{t})$.
    \STATE Update selected block:
    \[
    w_{t+1}^{(i_{t})}=w_{t}^{(i_{t})}-\eta_{t}\,g_{t}.
    \]
    \STATE Keep other blocks unchanged:
    \[
    w_{t+1}^{(b)}=w_{t}^{(b)}\quad\forall b\neq i_{t}.
    \]
\ENDFOR
\STATE \textbf{Output:} $w_{T}$ (or any iterate $w_{t}$).
\end{algorithmic}
\end{algorithm}

\section*{Dry Run Example}
Consider the scalar function $f(w)=(w-3)^{2}$, i.e. $d=1$ and a single block $B=1$.  
Take $w_{0}=0$, constant step size $\eta_{t}=0.1$, and run three iterations.

\begin{tabular}{c|c|c|c}
$t$ & $w_{t}$ & $\nabla f(w_{t})=2(w_{t}-3)$ & $f(w_{t})$\\ \hline
0 & $0$ & $-6$ & $9$\\
1 & $0-0.1(-6)=0.6$ & $2(0.6-3)=-4.8$ & $(0.6-3)^{2}=5.76$\\
2 & $0.6-0.1(-4.8)=1.08$ & $2(1.08-3)=-3.84$ & $(1.08-3)^{2}=3.6864$\\
3 & $1.08-0.1(-3.84)=1.464$ & $2(1.464-3)=-3.072$ & $(1.464-3)^{2}=2.3710$\\
\end{tabular}

The objective value decreases from $9$ to $2.37$ after three updates.

\newpage
\section*{Assumptions}
\begin{enumerate}[label=\textbf{A\arabic*:}, leftmargin=*]
\item \textbf{Blockwise $L$‑smoothness.}  
For each block $b$ there exists $L_{b}>0$ such that for all $u,v\in\mathbb{R}^{d}$,
\[
f\big(u^{(b)},v^{(-b)}\big)\le 
f\big(v^{(b)},v^{(-b)}\big)
+ \langle \nabla_{b}f(v),\,u^{(b)}-v^{(b)}\rangle
+\frac{L_{b}}{2}\|u^{(b)}-v^{(b)}\|^{2},
\tag{1}
\]
where $v^{(-b)}$ denotes all blocks except $b$.
\item \textbf{Lower boundedness.}  
$f$ has a finite infimum $f^{\inf}>-\infty$ and we denote $f^{*}= \inf_{w}f(w)$.
\item \textbf{Step‑size restriction.}  
The (possibly constant) step size satisfies
\[
0<\eta_{t}\le \frac{1}{L_{\max}},\qquad 
L_{\max}:=\max_{b}L_{b}. \tag{2}
\]
\item \textbf{Unbiased block selection.}  
The random block $i_{t}$ is independent of $w_{t}$ and satisfies  
$\mathbb{E}[\,\mathbf{e}_{i_{t}}\,]=\frac{1}{B}\sum_{b=1}^{B}\mathbf{e}_{b}$, where $\mathbf{e}_{b}$ is the indicator vector of block $b$.
\end{enumerate}

\section*{Main Convergence Theorem}
\begin{theorem}[Expected Gradient Norm Decay]\label{thm:main}
Assume \textbf{A1–A4} and let the step size be constant $\eta_{t}\equiv\eta\le 1/L_{\max}$.  
Define the random iterate $w_{\tau}$ where $\tau$ is drawn uniformly from $\{0,\dots ,T-1\}$ and independent of the algorithmic randomness. Then
\[
\boxed{
\mathbb{E}\!\left[\big\|\nabla_{i_{\tau}}f(w_{\tau})\big\|^{2}\right]
\;\le\;
\frac{2\,\big(f(w_{0})-f^{*}\big)}{\eta\,T}
}
\tag{3}
\]
Consequently the expected squared norm of a randomly selected block gradient converges to zero at rate $O(1/T)$.
\end{theorem}

\section*{Theoretical Properties}
\begin{enumerate}[label=\textbf{P\arabic*:}, leftmargin=*]
\item \textbf{Stability.} The requirement $\eta\le 1/L_{\max}$ guarantees that each block update is a descent step in expectation (see inequality (5) below). Larger step sizes may break (1) and cause divergence.
\item \textbf{Hyper‑parameter sensitivity.} The bound (3) scales linearly with $1/\eta$: a smaller step size yields a tighter bound but slower practical progress; a larger step size (still $\le 1/L_{\max}$) accelerates convergence at the cost of a larger constant factor.
\item \textbf{Comparison to full‑gradient SGD.} When $B=1$, Theorem~\ref{thm:main} reduces to the classic non‑convex SGD bound $O(1/(\eta T))$ for deterministic gradients. For $B>1$, the per‑iteration computational cost is reduced proportionally to $1/B$, while the convergence rate in terms of iterations remains the same.
\item \textbf{Special case – diminishing step size.} If $\eta_{t}=c/\sqrt{t+1}$ with $c\le 1/L_{\max}$, a standard telescoping argument yields $\min_{0\le t<T}\mathbb{E}\!\big[\|\nabla_{i_{t}}f(w_{t})\|^{2}\big]=O(1/\sqrt{T})$.
\end{enumerate}

\section*{Discussion}
The theorem shows that even without convexity, random block updates achieve the same $O(1/T)$ decay of the expected squared block gradient norm as full‑gradient descent, provided the blockwise smoothness constants are respected. This justifies the use of RBCD in large‑scale non‑convex optimisation where evaluating the full gradient is prohibitive.

\newpage
\section*{Preliminaries (Definitions and Lemmas)}
\begin{lemma}[Blockwise Descent Lemma]\label{lem:descent}
Under \textbf{A1} and for any block $b$,
\[
f\big(w - \eta \mathbf{e}_{b}\nabla_{b}f(w)\big)
\;\le\;
f(w) - \eta\big\|\nabla_{b}f(w)\big\|^{2}
+\frac{L_{b}\eta^{2}}{2}\big\|\nabla_{b}f(w)\big\|^{2}.
\tag{4}
\]
\end{lemma}
\begin{proof}
Apply the blockwise smoothness inequality (1) with  
$u^{(b)} = w^{(b)}-\eta\nabla_{b}f(w)$ and $v=w$.  
All other blocks are unchanged, i.e. $u^{(-b)}=v^{(-b)}=w^{(-b)}$.  
Then
\[
\begin{aligned}
f\big(w-\eta\mathbf{e}_{b}\nabla_{b}f(w)\big)
&\le f(w) + \langle \nabla_{b}f(w), -\eta\nabla_{b}f(w)\rangle
   +\frac{L_{b}}{2}\|\eta\nabla_{b}f(w)\|^{2}\\
&= f(w) -\eta\|\nabla_{b}f(w)\|^{2}
   +\frac{L_{b}\eta^{2}}{2}\|\nabla_{b}f(w)\|^{2}.
\end{aligned}
\]
\end{proof}

\section*{Main Proof (Step‑by‑Step Derivation)}
We prove Theorem~\ref{thm:main}. All expectations are taken w.r.t.\ the randomness of the block indices $\{i_{t}\}$.

\begin{enumerate}[label=\textbf{[Step \arabic*]}]
\item \textbf{Apply Lemma~\ref{lem:descent} to the selected block.}  
For the random block $i_{t}$ we have
\[
f(w_{t+1})
\le f(w_{t}) -\eta\big\|\nabla_{i_{t}}f(w_{t})\big\|^{2}
+ \frac{L_{i_{t}}\eta^{2}}{2}\big\|\nabla_{i_{t}}f(w_{t})\big\|^{2}.
\tag{5}
\]

\item \textbf{Take conditional expectation w.r.t.\ $i_{t}$.}  
Using uniform sampling,
\[
\mathbb{E}_{i_{t}}\!\big[\|\nabla_{i_{t}}f(w_{t})\|^{2}\big]
= \frac{1}{B}\sum_{b=1}^{B}\|\nabla_{b}f(w_{t})\|^{2},
\qquad
\mathbb{E}_{i_{t}}\!\big[L_{i_{t}}\|\nabla_{i_{t}}f(w_{t})\|^{2}\big]
= \frac{1}{B}\sum_{b=1}^{B}L_{b}\|\nabla_{b}f(w_{t})\|^{2}.
\tag{6}
\]

\item \textbf{Insert (6) into (5) and use $\eta\le 1/L_{\max}$.}  
Since $L_{b}\le L_{\max}$,
\[
\begin{aligned}
\mathbb{E}_{i_{t}}\!\big[f(w_{t+1})\big]
&\le f(w_{t}) -\eta\,\mathbb{E}_{i_{t}}\!\big[\|\nabla_{i_{t}}f(w_{t})\|^{2}\big]
   +\frac{\eta^{2}}{2}\mathbb{E}_{i_{t}}\!\big[L_{i_{t}}\|\nabla_{i_{t}}f(w_{t})\|^{2}\big]\\
&\le f(w_{t}) -\eta\,\mathbb{E}_{i_{t}}\!\big[\|\nabla_{i_{t}}f(w_{t})\|^{2}\big]
   +\frac{\eta^{2}L_{\max}}{2}\,\mathbb{E}_{i_{t}}\!\big[\|\nabla_{i_{t}}f(w_{t})\|^{2}\big]\\
&= f(w_{t}) -\Big(\eta-\frac{\eta^{2}L_{\max}}{2}\Big)
   \mathbb{E}_{i_{t}}\!\big[\|\nabla_{i_{t}}f(w_{t})\|^{2}\big].
\end{aligned}
\tag{7}
\]

\item \textbf{Simplify the coefficient using (2).}  
Because $\eta\le 1/L_{\max}$,
\[
\eta-\frac{\eta^{2}L_{\max}}{2}
\;\ge\;
\frac{\eta}{2}.
\tag{8}
\]

\item \textbf{Obtain a one‑step expected descent inequality.}  
Combining (7) and (8) gives
\[
\mathbb{E}_{i_{t}}\!\big[f(w_{t+1})\big]
\le f(w_{t}) -\frac{\eta}{2}\,
      \mathbb{E}_{i_{t}}\!\big[\|\nabla_{i_{t}}f(w_{t})\|^{2}\big].
\tag{9}
\]

\item \textbf{Take total expectation over the history up to time $t$.}  
Denote $\mathbb{E}_{t}$ the expectation conditioned on $\{i_{0},\dots ,i_{t-1}\}$.  
Since (9) holds conditionally,
\[
\mathbb{E}_{t}\!\big[f(w_{t+1})\big]
\le f(w_{t}) -\frac{\eta}{2}\,
      \mathbb{E}_{t}\!\big[\|\nabla_{i_{t}}f(w_{t})\|^{2}\big].
\tag{10}
\]

\item \textbf{Iterate (10) and telescope.}  
Taking expectation over all randomness and summing $t=0$ to $T-1$,
\[
\begin{aligned}
\mathbb{E}\!\big[f(w_{T})\big]
&\le f(w_{0}) -\frac{\eta}{2}\sum_{t=0}^{T-1}
      \mathbb{E}\!\big[\|\nabla_{i_{t}}f(w_{t})\|^{2}\big].
\end{aligned}
\tag{11}
\]

\item \textbf{Use lower boundedness \textbf{A2}.}  
Since $f(w_{T})\ge f^{*}$,
\[
f(w_{0})-f^{*}\;\ge\;
\frac{\eta}{2}\sum_{t=0}^{T-1}
      \mathbb{E}\!\big[\|\nabla_{i_{t}}f(w_{t})\|^{2}\big].
\tag{12}
\]

\item \textbf{Average over the $T$ iterates.}  
Divide (12) by $T$,
\[
\frac{1}{T}\sum_{t=0}^{T-1}
      \mathbb{E}\!\big[\|\nabla_{i_{t}}f(w_{t})\|^{2}\big]
\le \frac{2\big(f(w_{0})-f^{*}\big)}{\eta\,T}.
\tag{13}
\]

\item \textbf{Introduce the uniformly random index $\tau$.}  
Let $\tau$ be independent of the algorithm and uniformly drawn from $\{0,\dots ,T-1\}$.  
Then by definition of expectation,
\[
\mathbb{E}\!\big[\|\nabla_{i_{\tau}}f(w_{\tau})\|^{2}\big]
= \frac{1}{T}\sum_{t=0}^{T-1}
      \mathbb{E}\!\big[\|\nabla_{i_{t}}f(w_{t})\|^{2}\big].
\tag{14}
\]

\item \textbf{Combine (13) and (14) to obtain the claimed bound.}
\[
\boxed{
\mathbb{E}\!\big[\|\nabla_{i_{\tau}}f(w_{\tau})\|^{2}\big]
\le \frac{2\big(f(w_{0})-f^{*}\big)}{\eta\,T}
}
\tag{15}
\]
which is exactly inequality (3) of Theorem~\ref{thm:main}. ∎
\end{enumerate}

\section*{Rate Derivation (With Constants)}
From (15) we read off the explicit constants:
\[
C_{0}=2\big(f(w_{0})-f^{*}\big),\qquad
\text{rate } \; \mathbb{E}\!\big[\|\nabla_{i_{\tau}}f(w_{\tau})\|^{2}\big]\le \frac{C_{0}}{\eta\,T}.
\]
If a diminishing step size $\eta_{t}=c/\sqrt{t+1}$ is used, the same telescoping argument yields
\[
\min_{0\le t<T}\mathbb{E}\!\big[\|\nabla_{i_{t}}f(w_{t})\|^{2}\big]
\;\le\;
\frac{2\big(f(w_{0})-f^{*}\big)}{c}\,\frac{1}{\sqrt{T}}.
\]

\section*{Special Cases}
\begin{itemize}
\item \textbf{Single block ($B=1$).}  The algorithm reduces to standard (deterministic) gradient descent; Theorem~\ref{thm:main} recovers the classical $O(1/T)$ bound for non‑convex smooth functions.
\item \textbf{Identical Lipschitz constants $L_{b}=L$.}  Then $L_{\max}=L$ and the admissible step size simplifies to $\eta\le 1/L$, matching the usual smoothness condition for full‑gradient methods.
\item \textbf{Sparse updates.}  If only a small subset of blocks have non‑zero gradients, the bound (3) still holds because the expectation automatically averages over zero‑gradient blocks, potentially leading to faster practical convergence.
\end{itemize}

\end{document}